%------------------------------------------------------------------------------%
\section{Séries de Potências}

Nesse capítulo vamos utilizar a teoria desenvolvida para o estudo de séries numéricas
para construir as Séries de Potências, que podemos imaginar como polinômios de
grau infinito.
Para estudar a convergência das Séries de Potência, vamos considerar que para
cada valor fixo de $x$ temos uma série numérica que pode ser analisada com
os resultados do capítulo anterior. Porém,
estaremos também interessados em estudar essas séries como funções de $x$.
Na próxima seção apresentaremos as Séries de Taylor, que nos mostram
como construir séries de potência para funções que conhecemos,
desejamos estudar ou aproximar.
Começamos apresentando a definição da Série de Potências.

%------------------------------------------------------------------------------%
\begin{definicao}{Série de Potências}{def:Serie-potencia}
  \emph{Série de Potências}\index{Série!potências} centrada em $a$
  é uma série de funções da forma
  \[
  	\sum_{n=0}^\infty c_n(x-a)^n = c_0 + c_1(x-a) + c_2(x-a)^2
                                 + \cdots + c_n(x-a)^n + \cdots
 	\]
  onde $c_n$ e $a$ são constantes e $x\in\R$ é uma variável.
\end{definicao}
%------------------------------------------------------------------------------%

A seguir temos um exemplo de uma Série de Potências e a função que corresponde
a sua soma.

%------------------------------------------------------------------------------%
\begin{example}
  Determine a soma da série
  \(\;
    \sum_{n=0}^\infty x^n   \;.
  \)
  \solution
  Primeiro note que essa é uma Série de Potências centrada em $a=0$
  com coeficientes $c_n=1$, para todo $n=0, 1, 2,\dots$
  Consideramos agora o valor de $x$ fixo e observamos que essa é uma Série Geométrica
  com $r=x$, portanto convergente para todo $x$ tal que $\abs{x}<1$.
  Além disso, usando a fórmula da
  soma da Série Geométrica~\eqref{eq:soma-serie-geometrica},
  nos pontos onde a série converge, temos
  \[
  	\sum_{n=0}^\infty x^n = \sum_{n=1}^\infty x^{n-1} = \frac{1}{1-x}  \;.
 	\]
 	Como essa série diverge para todo $\abs{x}\geq 1$ essa igualdade só vale
  para o intervalo $(-1, 1)$.
\end{example}
%------------------------------------------------------------------------------%

\begin{figure}
  \begin{center}%
    \scalebox{0.9}{% Title: gl2ps_renderer figure
% Creator: GL2PS 1.4.2, (C) 1999-2020 C. Geuzaine
% For: Octave
% CreationDate: Tue Sep 28 14:08:49 2021
\setlength{\unitlength}{1pt}
\begin{picture}(0,0)
\includegraphics[width=17cm,height=7cm]{fig_taylor_1div1mx}
\end{picture}%
\begin{picture}(481,198)(0,0)
\fontsize{10}{0}\selectfont\put(43.3701,19.6115){\makebox(0,0)[t]{\textcolor[rgb]{0.15,0.15,0.15}{{-1}}}}
\fontsize{10}{0}\selectfont\put(428.882,19.6115){\makebox(0,0)[t]{\textcolor[rgb]{0.15,0.15,0.15}{{1}}}}
\fontsize{10}{0}\selectfont\put(231.13,62.2846){\makebox(0,0)[r]{\textcolor[rgb]{0.15,0.15,0.15}{{1}}}}
\fontsize{10}{0}\selectfont\put(231.13,97.477){\makebox(0,0)[r]{\textcolor[rgb]{0.15,0.15,0.15}{{2}}}}
\fontsize{10}{0}\selectfont\put(231.13,132.669){\makebox(0,0)[r]{\textcolor[rgb]{0.15,0.15,0.15}{{3}}}}
\fontsize{10}{0}\selectfont\put(231.13,167.862){\makebox(0,0)[r]{\textcolor[rgb]{0.15,0.15,0.15}{{4}}}}
\fontsize{10}{0}\selectfont\put(438.52,97.477){\makebox(0,0)[l]{\textcolor[rgb]{0,0,0}{{$P_1$}}}}
\fontsize{10}{0}\selectfont\put(438.52,132.669){\makebox(0,0)[l]{\textcolor[rgb]{0,0,0}{{$P_2$}}}}
\fontsize{10}{0}\selectfont\put(438.52,167.862){\makebox(0,0)[l]{\textcolor[rgb]{0,0,0}{{$P_3$}}}}
\fontsize{10}{0}\selectfont\put(438.52,196.016){\makebox(0,0)[l]{\textcolor[rgb]{0,0,0}{{$P_4$}}}}
\fontsize{10}{0}\selectfont\put(376.838,185.458){\makebox(0,0)[r]{\textcolor[rgb]{0,0,0}{{\(\dfrac{1}{1-x}\)}}}}
\end{picture}
}%
    \vspace{-\baselineskip}%
  \end{center}
  \caption{Comparação da função \(\sfrac{1}{(1-x)}\) com seus\\Polinômios de Taylor.}
  \label{fig:funcao_tayplor_1div1mx}
\end{figure}

A Figura~\ref{fig:funcao_tayplor_1div1mx} exibe o gráfico da função
\(
  \sfrac{1}{(1-x)}
\)
e as quatro primeiras somas parciais da série de potências.

%------------------------------------------------------------------------------%
Com base no resultado desse exemplo podemos usar a série truncada como uma
aproximação para a função, desde que $\abs{x}<1$, isso é,
\[
	f(x) = \frac{1}{1-x} \;\approx\; \sum_{n=0}^{N} x^n   \;.
\]
O próximo exemplo apresenta uma generalização desse resultado.

%------------------------------------------------------------------------------%
\begin{example}
  Para quaisquer constantes reais $a$ e $\beta\neq 0$, calcule a soma da série
  \[
    \sum_{n=0}^\infty \beta^n(x-a)^n   \;.
  \]

  \solution
  Essa é uma Série de Potências centrada em $a$ com coeficientes $c_n=\beta^n$
  que podemos escrever como
  \[
  	\sum_{n=0}^\infty \big(\beta(x-a)\big)^n   \;.
 	\]
  Para cada $x$ fixo, esta é uma série geométrica com $r=\beta(x-a)$.
  Ela converge se, e somente se, $\abs{r}<1$, ou seja, quando
  \begin{gather*}
  	\abs{\beta(x-a)}        \;<\; 1                                     \\[3mm]
  	-\frac{1}{\abs{\beta}}  \;<\;  x-a  \;<\; \frac{1}{\abs{\beta}}     \\[3mm]
  	a-\frac{1}{\abs{\beta}} \;<\;  x    \;<\; a + \frac{1}{\abs{\beta}}   \;.
  \end{gather*}
  Novamente usamos a fórmula da
  soma da Série Geométrica~\eqref{eq:soma-serie-geometrica}
  nos pontos onde a série converge
  \[
    \sum_{n=0}^\infty \beta^n(x-a)^n =
    \sum_{n=1}^\infty \big(\beta(x-a)\big)^{n-1} =
    \frac{1}{1-\beta(x-a)}   \;.
  \]
\end{example}
%------------------------------------------------------------------------------%

%------------------------------------------------------------------------------%
\begin{example}
  Calcule a soma da série
  \(\;
    \sum_{n=0}^\infty \left(\frac{x-2}{2}\right)^n   \;.
  \)
  \solution
  Essa série é um caso particular do exemplo anterior
  com $\beta=\sfrac{1}{2}$ e $a=2$, temos então que
  \[
    \sum_{n=0}^\infty \left(\frac{x-2}{2}\right)^n
    = \frac{1}{1-\frac{x-2}{2}}
    = \frac{2}{4-x}   \;,
  \]
  que converge para
  \begin{align*}
    a-\frac{1}{\abs{\beta}} <          & \; x < a + \frac{1}{\abs{\beta}}        \;,\\[3mm]
    2 - \abs{\frac{1}{\sfrac{1}{2}}} < & \; x < 2 + \abs{\frac{1}{\sfrac{1}{2}}} \;,\\[3mm]
    0 <                                & \; x < 4   \;.
  \end{align*}
\end{example}
%------------------------------------------------------------------------------%

\OnlineBox{https://colab.research.google.com/drive/19dbJqWuwaQGKtibaxspj0IQJyntL_UKo?usp=sharing}
{
  Código Python para ilustrar algumas séries de potências.
}

Os exemplos anteriores mostraram como podemos usar a definição de convergência de
séries numéricas para estudar a convergência das séries de potências.
Porém, existem resultados específicos que podemos utilizar para as Séries de Potência,
por exemplo, toda série de potências converge no seu centro, chamamos esse caso
de \emph{Convergência Trivial}\index{Convergência!trivial}. Podemos comprovar
isso simplesmente avaliando cada termo da série em $x=a$
\[
  \sum_{n=0}^\infty c_n(x-a)^n \,\evalat{x=a}{}
    = c_0 + c_1(a-a) + c_2(a-a)^2 + \cdots
    = c_0   \;.
\]
O teorema a seguir apresenta uma propriedade das séries de potências
que simplifica a analise de convergência.

%------------------------------------------------------------------------------%
\begin{teorema}{Convergência da Série de Potências}{teo:convergencia-serie-potencias}
  Se a Série de Potências\index{Série!potências!convergência}
  \[
  	\sum_{n=0}^\infty c_n(x-a)^n
 	\]
  converge em $x-a=L\neq 0$, então ela converge absolutamente para todo $x$ com
  \[
  	\abs{x-a}<\abs{L}   \;.
  \]
  Se a série diverge em $x-a=M$, então ela diverge para todo $x$ tal que
  \[
  	\abs{x-a}>\abs{M}   \;.
  \]
\end{teorema}
%------------------------------------------------------------------------------%
\begin{proof}
  \emph{Parte 1}

  Suponha que
  \[
  	\sum_{n=0}^\infty c_n L^n
 	\]
  converge, então, pelo Teste de Divergência~\ref{teo:teste-divergencia}
  \( \lim (c_n L^n) = 0 \).
  Assim, existe $N\in\N$ tal que
  \[
  	\abs{ c_n L^n } < 1 \qquad \forall n > N \;.
 	\]
 	Dividindo os dois lados por $\abs{L}^n$ e multiplicando por $\abs{x-a}^n$, temos
  \begin{equation}
    \label{eq:proof_desigualdade_L}
  	\abs{c_n}\abs{x-a}^n \;<\; \abs{\frac{x-a}{L}}^n \qquad \forall n>N   \;.
  \end{equation}
  Note que, se $x$ é tal que $\abs{x-a} < \abs{L}$, a série geométrica
  \[
    \sum_{n=0}^\infty \abs{\frac{x-a}{L}}^n
  \]
  converge.
  Então, pelo Teste da Comparação~\ref{teo:teste-comparacao}
  usando a desigualdade~\eqref{eq:proof_desigualdade_L}
  a série
  \[
  	\sum_{n=0}^\infty \abs{c_n}\abs{x-a}^n
  \]
  converge. Isso é o mesmo que dizer que a série
  \[
  	\sum_{n=0}^\infty c_n(x-a)^n
  \]
  converge absolutamente.

  \emph{Parte 2}

  Assuma que a série diverge para
  \[
    r - a = M \neq 0
  \]
  e suponha, por absurdo, que existe $s$ satisfazendo
  \[
    \abs{s-a} > \abs{M}
  \]
  onde a série converge.
  Usando a primeira parte do teorema sabemos que se a série converge em $s$
  ela deve convergir para todo $x$ tal que
  \[
    \abs{x-a} < \abs{s-a}
  \]
  Em particular deve convergir em $r$ o que é uma contradição.
  Provamos então que não pode existir nenhum ponto $s$ convergente
  tal que
  \[
    \abs{s-a} > \abs{M} \;.
  \]
\end{proof}
%------------------------------------------------------------------------------%

Esse teorema é útil para encontrar intervalos onde a série converge
absolutamente e onde a série diverge, o próximo exemplo emprega esse resultado.

%------------------------------------------------------------------------------%
\begin{example}
  Determine para quais valores de $x$ a série
  \[
    \sum_{n=0}^\infty (-1)^{n+1}\frac{(x-2)^n}{n}
  \]
  é convergente.
  \solution
  Note que quando $x-2=\pm 1$ podemos decidir com facilidade se a série converge ou não.
  Considerando o caso $x-2=1$ temos $x=3$ que gera a série numérica
  \[
    \sum_{n=0}^\infty\frac{ (-1)^{n+1}}{n}
  \]
  que converge pois é a Série Harmônica Alternada.
  Assim, pelo Teorema~\ref{teo:convergencia-serie-potencias},
  podemos afirmar que essa série converge absolutamente para todo $x$ tal que $\abs{x-2}<1$,
  ou seja, a série converge absolutamente no intervalo $(1, 3)$.
  Note que no ponto~$x=3$ a série converge condicionalmente,
  por se tratar da Série Harmônica Alternada.

  Tomando agora $x-2=-1$, isso é $x=1$, obtemos
  \[
    \sum_{n=0}^\infty\frac{\,(-1)^{n+1}(-1)^n\,}{n}
    = \sum_{n=0}^\infty\frac{-1\,}{n} \;.
  \]
  Como a Série Harmônica diverge, essa última série também diverge.
  Assim pelo Teorema~\ref{teo:convergencia-serie-potencias},
  a série diverge sempre que $\abs{x-a}>1$, ou seja,
  a série diverge em $(-\infty,\,1]\cup (3,\,\infty)$.
  O gráfico a seguir exibe a região de convergência dessa série.
  No interior do intervalo a convergência é absoluta e no ponto $x=3$ é condicional.
  %----------------------------------------------------------------------------%
  \IncludeTexGraphicCentred{fig_interval_exemplo_1}[-1.8]
  %----------------------------------------------------------------------------%
\end{example}
%------------------------------------------------------------------------------%

Note que o último teorema garante que os pontos onde uma série de potências
converge estão sempre em um intervalo, que chamamos de intervalo de convergência.

%------------------------------------------------------------------------------%
\begin{definicao}{Intervalo de convergência}{def:Intervalo-convergencia}
  O intervalo dos pontos onde uma Série de Potências converge é chamado de
  \emph{Intervalo de Convergência}%
  \index{Série!potências!intervalo de convergência}%
  \index{Intervalo!convergência}.
\end{definicao}
%------------------------------------------------------------------------------%

O teorema sobre a convergência das séries de potências pode ser
reescrito de uma forma mais poderosa, apresentada a seguir.

%------------------------------------------------------------------------------%
\begin{teorema}{Raio de Convergência}{teo:raio-convergencia}
  O Intervalo de Convergência de uma
  Série de Potências\index{Série!potências!convergência} centrada em $a$
  \[
  	\sum_{n=0}^\infty c_n(x-a)^n
 	\]
  sempre será um intervalo centrado em $a$ com raio $R$.
  Chamamos o número $R$ de
  \emph{Raio de Convergência}%
  \index{Série!potências!raio de convergência}%
  \index{Raio de convergência}
  da Série de Potências e ele se enquadra em um dos casos:
  \begin{description}[leftmargin=10mm, labelwidth=17mm, itemindent=20mm]
    \item[$R=0$]      Convergência trivial, a série converge apenas em $x=a$;
    \item[$R=\infty$] A série converge absolutamente para todo $x\in\R$;
    \item[$R>0$]      A série converge absolutamente para todo $x$ satisfazendo
      \[
        \abs{x-a}<R
      \]
      e diverge para todo $x$ tal que
      \[
       \abs{x-a}>R   \;.
      \]
  \end{description}
  \IncludeTexGraphicCentred{fig_interval_raio_convergencia}[-1.8][-1]
\end{teorema}
%------------------------------------------------------------------------------%
\begin{proof}
  Queremos mostrar que se não estivermos em um dos dois primeiros casos
  precisamos estar no terceiro, isso é, se o conjunto dos valores de $x$
  onde a série converge
  \[
    \Omega = \left\{ \; x\in\R \st \sum_{n=0}^\infty c_n(x-a)^n \text{ converge } \right\}
  \]
  não for $\Omega=\{a\}$ nem $\Omega=\R$ ele precisa ser o intervalo
  \[
    \Omega = [\, a-R, a+R \,]
  \]
  para algum número finito $R$.

  Como não estamos na situação em que a convergência ocorre para toda a reta,
  existe um ponto $u$ onde a série diverge,
  Assim, pelo Teorema~\ref{teo:convergencia-serie-potencias},
  temos que a série diverge para todo $x$ tal que $\abs{x-a}>d$, onde $d=\abs{u-a}$.
  Concluímos então que $\Omega$ é um conjunto limitado
  \[
  	\Omega\subset [\, a-d, a+d \,]   \;.
  \]
  Como também não estamos na situação de convergência trivial,
  existe $x\in \Omega$ diferente de $a$. Tomamos então $R$
  como o menor número maior ou igual a $\abs{x-a}$ para todo $x\in \Omega$, isso é,
  (veja a Definição~\ref{def:inf-sup})
  \begin{equation}
    \label{eq:r_sup}
    R = \sup_{x\in \Omega} \,\abs{x-a}   \;.
  \end{equation}
  Como $\Omega$ é limitado e existe $x\neq a$ em $\Omega$ sabemos que $R>0$ e é finito.
  Agora, vamos mostrar que a série converge para qualquer $x$ tal que $\abs{x-a}<R$.
  Como $R$ é tal que a condição~\eqref{eq:r_sup} vale, existe $u\in \Omega$ tal que
  \[
    \abs{x-a} < \abs{u-a} \leq R
  \]
  então, pelo Teorema~\ref{teo:convergencia-serie-potencias},
  a série converge absolutamente em~$x$.
  Por outro lado, se $v$ é tal que
  \[
    R < \abs{v-a}
  \]
  então a série tem que divergir em $v$, pois caso contrário $v\in \Omega$ o que
  contraria~\eqref{eq:r_sup}.
\end{proof}
%------------------------------------------------------------------------------%

Usando o Teorema~\ref{teo:convergencia-serie-potencias} ao identificarmos
o intervalo onde a série converge absolutamente, $\abs{x-a}<R$,
sabemos automaticamente que a série diverge para $\abs{x-a}>R$.
Entretanto, o teorema não diz nada sobre os extremos do intervalo
de convergência, $\abs{x-a}=R$, precisamos analisar individualmente esses pontos.
Para determinar $R$ podemos usar os testes de convergência, em especial o
Teste da Razão (Teorema~\ref{teo:teste-razao}) ou o
Teste da Raiz (Teorema~\ref{teo:teste-raiz}). Os próximos exemplos mostram como
determinar o intervalo de convergência de algumas séries.

%------------------------------------------------------------------------------%
\begin{example}
  Analise a convergência da série de potências
  \(
  	\sum_{n=0}^\infty \frac{(x-2)^n}{10^n}   \;.
 	\)
  \solution
  O termo geral dessa série é
  \[
  	a_n=\frac{(x-2)^n}{10^n}
 	\]
  então podemos aplicar o Teste da Raiz no módulo do termo geral, ou seja,
  \[
  	\sqrt[n]{\abs{a_n}}                     =
    \sqrt[n]{\abs{\frac{(x-2)^n}{10^n}}}    =
    \sqrt[n]{\abs{\frac{x-2}{10}}^n}        =
    \frac{\abs{x-2}}{10}                    \;\rightarrow\;
    \frac{\abs{x-2}}{10}
 	\]
  Então a série converge absolutamente desde que
  \[
  	\frac{\abs{x-2}}{10} < 1
 	\]
  ou seja, converge pare
  \(
  	\abs{x-2} < 10
 	\)
  e diverge sempre que
  \(
    \abs{x-2}>10
  \).
  Assim, o raio de convergência é $R=10$.

  Resta analisar o que ocorre nos extremos $x-2=\pm 10$.
  Se $x-2=10$ temos a série
  \[
    \sum_{n=0}^\infty \frac{(x-2)^n}{10^n} =
    \sum_{n=0}^\infty \frac{10^n}{10^n} =
    \sum_{n=0}^\infty 1
  \]
  que diverge pelo Teste da Divergência.
  No caso $x-2=-10$, obtemos a série
  \[
    \sum_{n=0}^\infty \frac{(x-2)^n}{10^n} =
    \sum_{n=0}^\infty \frac{(-10)^n}{10^n} =
    \sum_{n=0}^\infty (-1)^n
  \]
  que também diverge pelo Teste da Divergência.

  Em resumo a série converge absolutamente no intervalo $(-8,\,12)$
  e diverge em $(-\infty,\,-8]\cup[12,\,\infty)$.
  O gráfico a seguir ilustra os intervalo de convergência da série.
  %---------------------------------------------------------------------------%
  \IncludeTexGraphicCentred{fig_interval_exemplo_2}[-1.5]
  %---------------------------------------------------------------------------%
\end{example}
%------------------------------------------------------------------------------%

\OnlineBox{https://colab.research.google.com/drive/1EugBhApxYNwAiXASe-HWgRLVs6EUePL2?usp=sharing}
{
  Código Python para ilustrar o raio de convergência das séries de potências.
}

%------------------------------------------------------------------------------%
\begin{example}
  Analise a convergência da série de potências
  \(
  	\sum_{n=0}^\infty (-1)^n\frac{x^n}{n!}  \;.
 	\)
  \solution
  Aplicando o Teste da Razão para encontrar o raio de convergência obtemos,
  \[
      \abs{\frac{a_{n+1}}{a_n}} =
      \abs{\frac{x^{n+1}}{(n+1)!}\frac{n!}{x^n}} =
      \abs{\frac{x}{(n+1)}}
      \;\rightarrow\; 0  \;.
  \]
  Como o limite é menor do que 1 para qualquer valor de $x$
  concluímos que a série converge em toda a reta real,
  ou seja, o raio de convergência da série é~$R=\infty$.
\end{example}
%------------------------------------------------------------------------------%

%------------------------------------------------------------------------------%
\begin{example}
  Analise a convergência da série de potências
  \(
  	\sum_{n=0}^\infty \abs{n^n x^n}  \;.
 	\)
  \solution
  Aplicando o Teste da Raiz para a série obtemos, para todo $x\neq 0$, que
  \[
  	\sqrt[n]{\abs{n^nx^n}}=n\abs{x} \;\rightarrow\; \infty  \;.
 	\]
  Então essa série de potências converge apenas em $x=0$,
  isto é, seu raio de convergência é $R=0$.
\end{example}
%------------------------------------------------------------------------------%

Resumindo, quando queremos testar a convergência de uma Série de Potências,
normalmente empregamos os seguintes passos:
\begin{enumerate}[nosep]
  \item Usar o teste da razão, ou da raiz, para encontrar o intervalo onde
        a série converge absolutamente
        \[
          \abs{x-a} < R \qquad \text{ ou } \qquad a-R < x < a+R
        \]
  \item Se o intervalo de convergência absoluta for finito precisamos
        testar a convergência em cada extremo. Nesses pontos os testes
        da razão ou raiz provavelmente serão inconclusivos, portanto
        devemos usar o teste da comparação, integral ou série alternada.
\end{enumerate}
Observe que se o intervalo de convergência absoluta for finito, \(\abs{x-a} < R\),
podemos afirmar diretamente que a série diverge para \(\abs{x-a} > R\).

%------------------------------------------------------------------------------%
